\documentclass[12pt,a4paper]{article}

\usepackage[margin=1in]{geometry}
\usepackage{graphicx}
\usepackage{booktabs}
\usepackage{array}
\usepackage{longtable}
\usepackage{hyperref}
\usepackage{xcolor}
\usepackage{float}
\usepackage{amsmath}
\usepackage{enumitem}
\usepackage{caption}

\hypersetup{
    colorlinks=true,
    linkcolor=blue,
    urlcolor=blue,
    citecolor=blue
}

\setlength{\parskip}{0.5em}

\begin{document}

% ============================================================
% TITLE PAGE
% ============================================================

\begin{titlepage}
    \centering
    
    \vspace*{1.5cm}
    
    {\Huge \textbf{AI-Based Village Pond Planning System}\par}
    
    \vspace{0.4cm}
    
    {\Large Assignment 1 -- Phase 2\par}
    
    \vspace{1.5cm}
    
    {\large \textbf{Backend API and Terrain-Based Pond Analysis}\par}
    
    \vfill
    
    \begin{tabular}{rl}
        \textbf{Submitted By:} & Jaswinder Singh \\
        \textbf{Roll Number:} & \textit{Add Your Roll Number} \\
        \textbf{Department:} & Computer Science and Engineering \\
        \textbf{Institute:} & Indian Institute of Technology Bhilai \\
        \textbf{Assignment:} & Assignment 1 -- Phase 2 \\
        \textbf{Submission Date:} & 01/09/2026 \\
    \end{tabular}
    
    \vfill
    
    {\large Indian Institute of Technology Bhilai\par}
    
\end{titlepage}

% ============================================================
% TABLE OF CONTENTS
% ============================================================

\tableofcontents
\newpage

% ============================================================
\section{Introduction}

Planning a village pond requires identifying locations where rainfall runoff can be
effectively collected while considering terrain characteristics, drainage patterns,
catchment area, and environmental suitability.

The AI-Based Village Pond Planning System is designed as a geographical decision
support platform for analyzing terrain and identifying suitable locations for pond
construction. The system accepts contour data or a geographical location as input
and performs terrain-based analysis using multiple computational modules.

The backend is implemented as a REST API using FastAPI. The primary contour analysis
endpoint accepts an uploaded contour file and performs terrain analysis, candidate
generation, rainfall analysis, and multi-factor suitability evaluation.

The frontend presents the generated results through an interactive geographical
dashboard showing terrain information, candidate locations, catchment areas, and
suitability scores.

The implementation is designed as an extensible system so that additional
hydrological, environmental, and geographical factors can be integrated in future
phases.

% ============================================================
\section{Objectives}

The main objectives of the system are:

\begin{enumerate}
    \item Develop a backend API for terrain-based pond planning.
    \item Accept contour data as an input for geographical analysis.
    \item Extract terrain and elevation information from contour data.
    \item Generate terrain characteristics required for pond-site analysis.
    \item Analyze drainage and flow accumulation patterns.
    \item Identify multiple candidate locations for possible pond construction.
    \item Integrate rainfall-related information into the analysis.
    \item Perform multi-factor suitability analysis.
    \item Estimate catchment-related information for candidate locations.
    \item Return the complete analysis in structured JSON format.
    \item Provide interactive API documentation using Swagger/OpenAPI.
    \item Present the analysis results through an interactive frontend interface.
\end{enumerate}

% ============================================================
\section{Functional Requirements}

The system implements the following major functional requirements.

\begin{table}[H]
\centering
\caption{Functional Requirements}
\begin{tabular}{p{4cm} p{9cm}}
\toprule
\textbf{Requirement} & \textbf{Implementation} \\
\midrule
Contour Upload & Accepts contour files for terrain-based analysis. \\
Terrain Analysis & Extracts elevation and terrain information from contour data. \\
Candidate Generation & Identifies multiple possible locations for pond construction. \\
Rainfall Analysis & Includes rainfall-related information in the planning process. \\
Hydrological Analysis & Uses terrain and drainage characteristics for runoff concentration. \\
Suitability Analysis & Ranks candidate locations using geographical and environmental factors. \\
Catchment Analysis & Estimates the contributing catchment area for candidate locations. \\
Structured Response & Returns analysis results as structured JSON data. \\
API Documentation & Provides interactive Swagger/OpenAPI documentation through FastAPI. \\
Frontend Visualization & Displays candidates, catchment areas, and rankings on an interactive map. \\
Health Monitoring & Provides a health endpoint to verify backend service status. \\
\bottomrule
\end{tabular}
\end{table}

% ============================================================
\section{System Architecture}

The system follows a client-server architecture. The frontend provides an interactive
interface for selecting a geographical location or uploading contour information. The
request is sent to the FastAPI backend, which coordinates the terrain, hydrology,
rainfall, and suitability analysis modules.

The major processing stages are:

\begin{enumerate}
    \item \textbf{Input Selection}: User selects a geographical location or uploads contour data.
    \item \textbf{Contour Processing}: Backend extracts contour and elevation information.
    \item \textbf{Terrain Analysis}: Terrain characteristics such as elevation and resolution are generated.
    \item \textbf{Hydrological Processing}: Flow direction, drainage, and runoff characteristics are analyzed.
    \item \textbf{Rainfall Analysis}: Rainfall information is incorporated into the planning process.
    \item \textbf{Candidate Generation}: Potential pond locations are generated and evaluated.
    \item \textbf{Multi-Factor Suitability Analysis}: Multiple factors are combined for suitability scoring.
    \item \textbf{Catchment Analysis}: Catchment-related information is calculated for candidates.
    \item \textbf{JSON Response}: Backend returns the complete analysis as structured JSON.
    \item \textbf{Frontend Visualization}: Results are displayed on an interactive geographical dashboard.
\end{enumerate}

% ============================================================
\section{Technology Stack}

The system uses the following technologies.

\begin{table}[H]
\centering
\caption{Technology Stack}
\begin{tabular}{p{4cm} p{9cm}}
\toprule
\textbf{Technology} & \textbf{Purpose} \\
\midrule
Python & Primary programming language for backend development. \\
FastAPI & REST API framework for implementing backend endpoints. \\
Uvicorn & ASGI server to run the FastAPI application. \\
Pydantic & Request validation and structured response models. \\
NumPy & Numerical and terrain-related computations. \\
OpenCV & Image and contour-related processing. \\
Geospatial Services & Geographical, terrain, and location-related analysis. \\
Swagger/OpenAPI & Interactive API documentation and endpoint testing. \\
React / Modern Frontend & Interactive user interface for geographical visualization. \\
Leaflet / Map Visualization & Interactive display of terrain, candidates, and catchments. \\
Git/GitHub & Version control and source code hosting. \\
\bottomrule
\end{tabular}
\end{table}

% ============================================================
\section{API Design}

The backend provides multiple endpoints for the pond planning workflow.

\subsection{Primary Contour Analysis Endpoint}

\begin{center}
\texttt{POST /api/analyzeContour} \\
\texttt{http://localhost:8000/api/analyzeContour}
\end{center}

This endpoint accepts a contour file and performs the analysis pipeline. The endpoint
supports parameters related to terrain resolution, candidate generation, rainfall analysis,
runoff coefficient, and pond dimensions.

\subsection{Automatic Location Analysis Endpoint}

\begin{center}
\texttt{POST /api/analyzeLocation} \\
\texttt{http://localhost:8000/api/analyzeLocation}
\end{center}

This endpoint can be used for automatic analysis based on a location input.

\subsection{Health Endpoint}

\begin{center}
\texttt{GET /health} \\
\texttt{http://localhost:8000/health}
\end{center}

This endpoint verifies whether the FastAPI backend is running correctly.

\subsection{Root Endpoint}

\begin{center}
\texttt{GET /} \\
\texttt{http://localhost:8000/}
\end{center}

This serves the main index.html file containing the interactive frontend interface.

\subsection{API Documentation}

\begin{center}
\texttt{http://localhost:8000/docs}
\end{center}

The documentation contains the available endpoints, request parameters, request
schemas, and response schemas.

% ============================================================
\section{API Endpoint Parameters}

The primary contour analysis endpoint is \texttt{POST /api/analyzeContour} and accepts
multipart form data.

\begin{table}[H]
\centering
\caption{Contour Analysis Parameters}
\begin{tabular}{p{4cm} p{3cm} p{6cm}}
\toprule
\textbf{Parameter} & \textbf{Type} & \textbf{Description} \\
\midrule
file & File & Required contour file (KML/KMZ) used for terrain analysis. \\
resolution\_m & Number & Spatial resolution (2-100 m) for terrain analysis. Default: 10.0. \\
max\_candidates & Integer & Maximum number of candidate locations to generate (1-50). Default: 20. \\
rainfall\_years & Integer & Number of rainfall years considered (1-30). Default: 5. \\
runoff\_coefficient & Number & Coefficient for runoff estimation (0.0-1.0). Default: 0.30. \\
pond\_radius\_m & Number & Approximate radius for pond planning (10-150 m). Default: 40.0. \\
max\_pond\_depth\_m & Number & Maximum pond depth (1.0-6.0 m). Default: 3.0. \\
\bottomrule
\end{tabular}
\end{table}

\begin{table}[H]
\centering
\caption{Location Analysis Parameters}
\begin{tabular}{p{4cm} p{3cm} p{6cm}}
\toprule
\textbf{Parameter} & \textbf{Type} & \textbf{Description} \\
\midrule
location\_name & String & Required location name for geographical analysis. \\
analysis\_radius\_m & Number & Analysis radius in metres (500-10000 m). Default: 3000.0. \\
contour\_interval\_m & Number & Contour interval for DEM generation (1-20 m). Default: 5.0. \\
resolution\_m & Number & Spatial resolution (2-100 m). Default: 30.0. \\
max\_candidates & Integer & Maximum number of candidates (1-50). Default: 20. \\
rainfall\_years & Integer & Number of rainfall years (1-30). Default: 5. \\
runoff\_coefficient & Number & Runoff coefficient (0.0-1.0). Default: 0.30. \\
pond\_radius\_m & Number & Pond radius (10-150 m). Default: 40.0. \\
max\_pond\_depth\_m & Number & Maximum pond depth (1.0-6.0 m). Default: 3.0. \\
\bottomrule
\end{tabular}
\end{table}


% ============================================================
\section{Swagger API Documentation}

The FastAPI backend provides automatically generated Swagger documentation. The Swagger
interface displays the available API routes and allows the developer to inspect request
parameters and response schemas.

The contour analysis endpoint displayed in Swagger is \texttt{POST /api/analyzeContour}.
The request body uses \texttt{multipart/form-data}. The required input includes the
contour file, while additional optional parameters allow configuration of the terrain
and pond analysis.

\begin{figure}[H]
\centering
\includegraphics[width=\textwidth]{images/analyze_contour_endpoint.png}
\caption{Swagger documentation showing the Analyze Contour endpoint and request parameters}
\label{fig:swagger_endpoint}
\end{figure}

% ============================================================
\section{API Response Structure}

The contour analysis endpoint returns a structured JSON response containing:

\begin{itemize}
    \item Analysis status and input filename
    \item Terrain information including contour line count
    \item Minimum and maximum elevation
    \item Terrain resolution and grid dimensions
    \item Rainfall information
    \item Candidate generation information
    \item Land filtering information
    \item Suitability analysis
    \item Catchment-related information
    \item Error information when applicable
\end{itemize}

The Swagger response schema confirms that the API follows a structured response format
for integration with the frontend.

\begin{figure}[H]
\centering
\includegraphics[width=\textwidth]{images/analyze_contour_response.png}
\caption{Swagger response schema for the contour analysis endpoint}
\label{fig:swagger_response}
\end{figure}

% ============================================================
\section{Terrain and Hydrological Analysis}

The terrain analysis pipeline processes the contour information to obtain a
representation of the geographical area. The major terrain-related information includes
contour line information, elevation range, terrain resolution, grid dimensions, surface
slope, and drainage characteristics.

The generated terrain information is subsequently used by the hydrological analysis
modules. Hydrological processing is important because pond locations should preferably
be associated with areas where surface runoff can naturally accumulate. The analysis
considers terrain flow and drainage characteristics to identify potential areas for
water collection.

% ============================================================
\section{Candidate Generation and Suitability Analysis}

Potential pond locations are generated from the geographical analysis. Each candidate
location is evaluated using multiple factors. The suitability analysis combines terrain
and hydrological information with other available planning factors.

The frontend interface displays ranked candidate locations. The highest-ranked location
is identified as the top candidate, while additional candidate locations are also
displayed for comparison.

The suitability evaluation may include factors such as terrain characteristics, flow
accumulation, drainage properties, catchment area, rainfall information, land-related
constraints, accessibility, and environmental considerations.

% ============================================================
\section{Catchment Analysis}

A catchment represents the geographical area from which water contributes toward a
particular drainage outlet or candidate pond location. The system identifies
catchment-related information using terrain and drainage characteristics.

The catchment information is useful because a potential pond location should be
associated with sufficient upstream runoff. The frontend visualizes the catchment
area associated with the selected candidate location, allowing the user to visually
understand the selected pond candidate, the surrounding terrain, the drainage region,
the estimated catchment area, and other ranked candidate locations.

% ============================================================
\section{Frontend Visualization}

The project includes an interactive frontend dashboard for visualizing the pond
planning analysis. The frontend provides a geographical decision support interface
displaying terrain analysis information, candidate ranking, suitability scores, flow
accumulation information, valley position, drainage information, multi-factor
suitability, elevation, catchment area, storage-related estimates, an interactive
geographical map, candidate markers, catchment boundaries, and ranked pond locations.

\begin{figure}[H]
\centering
\includegraphics[width=\textwidth]{images/frontend_analysis_result.png}
\caption{Frontend interface showing terrain analysis, ranked pond candidates, catchment area, and geographical visualization}
\label{fig:frontend_result}
\end{figure}

% ============================================================
\section{Backend Health Testing}

The backend was tested using the health endpoint. An initial request was made to
\texttt{http://localhost:8000/api/health}, which returned:

\begin{verbatim}
{"detail":"Not Found"}
\end{verbatim}

This occurred because the correct health route is not under the \texttt{/api} prefix.
The correct endpoint is \texttt{http://localhost:8000/health}. The following command
was executed:

\begin{verbatim}
curl http://localhost:8000/health
\end{verbatim}

The backend returned:

\begin{verbatim}
{"status":"healthy"}
\end{verbatim}

This confirms that the FastAPI backend was running successfully.

\begin{figure}[H]
\centering
\includegraphics[width=\textwidth]{images/health_response.png}
\caption{Terminal output showing incorrect API route and successful health endpoint response}
\label{fig:health_response}
\end{figure}

% ============================================================
\section{Testing and Validation}

The backend and frontend were tested using the following validation steps:

\begin{enumerate}
    \item Verified that the FastAPI backend starts successfully.
    \item Opened the Swagger/OpenAPI documentation.
    \item Verified the availability of the contour analysis endpoint.
    \item Verified the request parameters for contour analysis.
    \item Inspected the API response schema.
    \item Tested the backend health endpoint.
    \item Verified that the correct health endpoint returns a healthy response.
    \item Verified that the frontend interface loads successfully.
    \item Verified that candidate locations are displayed on the geographical map.
    \item Verified that catchment information is visualized.
    \item Verified that suitability rankings are displayed.
\end{enumerate}

The successful health response obtained from the backend was:
\begin{verbatim}
{"status":"healthy"}
\end{verbatim}
This confirms that the backend API service was operational during testing.

% ============================================================
\section{API Endpoints Summary}

The following API endpoints were identified and used in the project.

\begin{table}[H]
\centering
\caption{API Endpoints}
\begin{tabular}{p{2cm} p{5cm} p{8cm}}
\toprule
\textbf{Method} & \textbf{Endpoint} & \textbf{Description} \\
\midrule
GET & / & Serves the frontend dashboard interface (index.html). \\
GET & /health & Checks whether the FastAPI backend is running correctly. \\
GET & /docs & Provides interactive Swagger/OpenAPI API documentation. \\
POST & /api/analyzeContour & Analyzes an uploaded contour file and performs terrain and pond planning analysis. \\
POST & /api/analyzeLocation & Performs location-based analysis and generates planning information. \\
\bottomrule
\end{tabular}
\end{table}

The primary API endpoint used for contour analysis in this assignment is:

\begin{center}
\fbox{\texttt{POST http://localhost:8000/api/analyzeContour}}
\end{center}

The location-based analysis endpoint is:

\begin{center}
\fbox{\texttt{POST http://localhost:8000/api/analyzeLocation}}
\end{center}

The health endpoint successfully tested is:

\begin{center}
\fbox{\texttt{GET http://localhost:8000/health}}
\end{center}

The root endpoint serving the frontend is:

\begin{center}
\fbox{\texttt{GET http://localhost:8000/}}
\end{center}

The Swagger documentation is available at:

\begin{center}
\fbox{\texttt{http://localhost:8000/docs}}
\end{center}

% ============================================================
\section{GitHub Repository}

The complete source code for the project is maintained in the following GitHub repository:

\begin{center}
\url{https://github.com/Jas-thinks/AI-POND-RECOMMENDATION}
\end{center}

The repository contains the backend implementation, terrain analysis modules, hydrological
analysis components, frontend application, and supporting project files.

% ============================================================
\section{Limitations}

The current implementation has several limitations.

\begin{enumerate}
    \item The accuracy of terrain analysis depends on the quality and density of the provided contour information.
    \item Geographical analysis can require significant computation for large areas.
    \item Suitability scores represent computational estimates and should not be considered final engineering approval.
    \item Catchment and storage estimates depend on the assumptions used by the terrain and hydrological models.
    \item Real-world pond construction requires additional field verification.
    \item Environmental, geological, and social factors may require further integration for production-level planning.
\end{enumerate}

% ============================================================
\section{Future Extensions}

The project has been designed to support further development. Possible future extensions include:

\begin{itemize}
    \item Integration with real-time rainfall datasets.
    \item Improved Digital Elevation Model generation.
    \item Advanced hydrological simulation.
    \item Soil-type analysis.
    \item Land-use and land-cover analysis.
    \item Satellite imagery integration.
    \item Environmental constraint analysis.
    \item Improved storage capacity estimation.
    \item Machine learning-based pond suitability prediction.
    \item Historical rainfall trend analysis.
    \item Persistent storage of analysis results.
    \item User authentication.
    \item Cloud deployment.
    \item Large-scale regional pond planning.
    \item Improved frontend visualization and reporting.
\end{itemize}

% ============================================================
\section{Conclusion}

The AI-Based Village Pond Planning System provides a backend-driven geographical
decision support system for identifying potential pond construction locations. The
system accepts contour or location-based input and performs terrain analysis,
hydrological processing, candidate generation, rainfall analysis, catchment-related
evaluation, and multi-factor suitability scoring.

The FastAPI backend exposes structured REST API endpoints for performing the analysis.
The primary contour analysis endpoint used in the system is:

\begin{center}
\texttt{POST http://localhost:8000/api/analyzeContour}
\end{center}

The backend was successfully verified using the health endpoint:

\begin{center}
\texttt{GET http://localhost:8000/health}
\end{center}

which returned:

\begin{verbatim}
{"status":"healthy"}
\end{verbatim}

Swagger/OpenAPI documentation was also successfully generated and accessed through:

\begin{center}
\texttt{http://localhost:8000/docs}
\end{center}

The frontend demonstrates the results through an interactive geographical dashboard
containing ranked pond candidates, terrain information, suitability scores, and
catchment visualization.

The project therefore provides a strong foundation for further development into a
comprehensive AI-assisted geographical planning platform for village pond construction
and rainwater harvesting.

\end{document}
